\section{Introduction}\label{sec:introduction}

Hypermetric inequalities connect cut polytopes, quadratic optimization,
and the geometry of lattices.  They are valid for every cut, and their
Boolean quadric images strengthen the usual linear and semidefinite
relaxations of binary quadratic problems.  We study the exact decision
problem of whether a rational point violates any hypermetric
inequality.  We prove that this problem is strongly NP-complete.
The hardness persists at points satisfying all triangle and perimeter
inequalities, whose elliptope matrix is positive definite, and which
satisfy every rounded positive semidefinite inequality up to any
prescribed fixed gonality cutoff.

The proof also determines the separation complexity of several related
families.  These include rounded positive semidefinite and gap-$1$
inequalities in cut form, Boros--Hammer inequalities in Boolean quadric
form, the Boolean quadric clique and cut inequalities, pure hypermetric
inequalities, and odd clique inequalities.  For the last family,
hardness holds even when only the unswitched inequalities with all
nonzero coefficients equal to $1$ are available.  These conclusions do
not follow merely from inclusion among inequality families.  We
construct points for which every possible violation in the larger
families encodes the same exact cover, and use switching and a small
orthogonal extension to obtain the families with coefficients in
$\{0,\pm1\}$.

\subsection{Relation to prior work}

\citet{AvisGrishukhin1993} formulated the recognition problem for the
hypermetric cone, established its membership in $\coNP$, and related it
to finding lattice points inside a sphere.  They also proved hardness
for problems in which a bound on the gonality of the inequality is part
of the input.  \citet{Avis2003} extended this complexity analysis to
negative-type, $k$-gonal, and gap inequalities.  These results leave a
logical distinction: restricting the coefficient norm can be hard even
when separation over the full class is easy.  Negative-type
inequalities illustrate the distinction, because full separation is a
positive semidefiniteness test.  The results for a prescribed gonality
therefore do not establish hardness of unrestricted hypermetric
separation.  \citet{DezaGrishukhin1997} also treated the equivalent
hypermetric correlation recognition problem and described its
NP-hardness as unknown.

The Boolean quadric families studied here have a separate polyhedral
history.  \citet{Padberg1989} developed the clique and cut inequalities,
and \citet{BorosHammer1993} studied the more general integer-coefficient
quadratic inequalities now bearing their names.  The covariance map
relates these families to cut inequalities; in particular, the
Boros--Hammer class is the Boolean quadric form of the rounded positive
semidefinite class.  \citet{LetchfordSorensen2012} established
separation reductions among these classes and hypermetric correlation
inequalities.  \citet{GKL2011,GKL2012} placed these families within the
theory of gap inequalities, originating in
\citet{LaurentPoljak1996}, and described the remaining complexity
questions.  \citet{Letchford2022} again identified hypermetric
separation as an open problem and recorded unknown separation
complexity for the clique, cut, and related Boolean quadric families.
The NP-completeness of testing $\gamma(b)=0$ for a specified vector
$b$ was already known \citep[\S2]{GKL2012}.  It does not determine
gap-$0$ separation, which asks whether a given point has any violating
coefficient vector with zero gap.

Two secondary accounts require a qualification of this history.
\citet[p.~23]{Liers2004} states that hypermetric separation is NP-hard
and cites Chapter~28 of \citet{DezaLaurent1997}.
\citet[\S6, ``The hypermetric inequalities'']{KrishnanTerlaky2005}
make an analogous assertion for the gap-$1$ class, which they call
hypermetric, citing \citet{AvisGrishukhin1993,Avis2003}.
Neither discussion supplies a reduction.  The primary complexity
papers and the later survey explicitly leave the unrestricted
question open.  We therefore distinguish those assertions from a
complete proof of the exact problem and its relaxation promises.

Hardness of separating split cuts in mixed-integer linear programming
\citep{CapraraLetchford2003} does not settle these questions.  Here the
input is a quadratic moment matrix, and in the binary case its
diagonal entries must equal its first-row entries.  At positive
definite moment points, these identities place the origin and the
lattice generators on a common sphere in the hypermetric representation.
A reduction for an unrestricted integer quadratic form
need not preserve them.  The construction below uses the fact that
every $0$--$1$ vector lies on the sphere through the origin centered at
$\one/2$.  It couples this sphere with exact-cover incidence vectors
and controls all integer points, including coefficient vectors of
arbitrarily large norm.

Conversely, the binary hard points are admissible general integer
moment matrices.  The same reduction therefore proves strong
NP-completeness of unrestricted integer-quadratic split separation,
whose complexity was left open by \citet{BuchheimTraversi2015}.
This consequence requires no separate integer construction.

Recent QCQP work also searches for Boros--Hammer inequalities by
solving nonconvex integer quadratic subproblems with variable
coefficients restricted to $[-2,2]$, in computational evaluations of
semidefinite relaxations \citep[\S4.1.2]{DeyEtAl2026}.

For algorithmic context, integer-quadratic separation methods can
exploit low rank and lattice structure \citep{BuchheimTraversi2015}.
Our algorithmic results give explicit reductions to exact
closest-vector problems in rational lattices.  Existing
closest-vector algorithms then yield bounds in the binary input
model.  We distinguish rank from support because they lead to
different parameter dependence.  We do not claim a new
closest-vector algorithm.

\subsection{Contributions and their scope}

The main contributions are as follows.

\begin{enumerate}
\item We give an exact-cover Gram construction whose violated
hypermetric inequalities correspond exactly to exact covers
(\cref{thm:exact-cover}).  Its negative slacks are all $-1/2$ in the
correlation representation.  A polynomially encoded scaling places
the distance point in $\MET(V)$ and its matrix $Z$ in the interior of
$\ELL(V)$, while excluding
every rounded positive semidefinite violation with coefficient sum of
absolute value at least $3$ (\cref{lem:scaling,thm:relaxation-hardness}).
This proves strong NP-completeness of unrestricted hypermetric
separation (\cref{cor:hypermetric-hardness}).

\item We transfer this construction to the gap-$1$, rounded positive
semidefinite, Boros--Hammer, and signed-coefficient Boolean quadric families
(\cref{thm:binary-families,thm:finite-families}).  A further orthogonal
extension proves strong NP-completeness for pure hypermetric and odd
clique separation, including unswitched positive odd clique
inequalities (\cref{thm:pure-clique-hardness}).  The witnesses for the
clique, cut, pure hypermetric, and odd clique families define facets;
we include a direct proof
(\cref{lem:clique-facets}).

\item We provide polynomial-size certificates for the infinite
families, an exact algorithm parameterized by rank, and an explicit
enumeration bound for separation above a normalized violation
threshold (\cref{sec:algorithms}).  The threshold bound is an XP bound:
its exponent depends on the reciprocal threshold.  A separate
general-integer example shows why restricting the allowed coefficient
vectors can destroy the tractability of low-rank separation.

\item We prove NP-completeness of gap-$0$ separation, with hard
points already in the metric polytope (\cref{thm:gap-zero-hardness}).
These points are outside the elliptope.  Gap-$0$ inequalities cannot
cut a positive semidefinite point, so this result has a different
algorithmic role from the main theorem.  Its reduction uses
PARTITION; we make no strong-hardness claim for this family.
\end{enumerate}

To the best of our knowledge, no prior work gives complete proofs of
the unrestricted hardness results in the first, second, and fourth
items with the stated domains and promises.
The novelty is the exact decision complexity of the stated families
and the explicit relaxation promises, rather than their definitions,
the covariance and switching operations, or the use of closest-vector
methods.  The earlier open-problem statements and separation
equivalences support this qualified claim; the secondary assertions
described above are part of its historical context.

These are exact separation results.  They rule out polynomial-time
exact oracles for the hard families unless $\mathrm P=\NP$, even after
the standard relaxations and any fixed gonality cutoff have been
enforced.  They do not compare cut-selection heuristics or establish
hardness at a fixed additive tolerance.  The constructed violations have inverse
polynomial size, so hardness is not caused by exponentially small
numerical slacks.  Their Euclidean-normalized size can nevertheless
decrease with the instance dimension; this is consistent with the
threshold algorithm in \cref{sec:algorithms}.

The paper first fixes the representations and separation conventions
(\cref{sec:preliminaries}).  \Cref{sec:core} develops the exact-cover
construction and its scaling.  \Cref{sec:classes} treats the related
binary families.  \Cref{sec:algorithms} gives certificates and
algorithms, and \cref{sec:gap-zero} treats gap-$0$ separation.
\Cref{sec:discussion} discusses the implications and the remaining
question for general gap inequalities at positive semidefinite
points.
